\section{地震・弾性波探査}

\subsection{平成23年度 問29}
\label{q:h23-reg-applied-29}

\noindent\textbf{原典：} sources/H23/17\_applied\_science.pdf、PDF物理10頁。

\textbf{公式正答：} 5（single）。

\paragraph{問題の要旨} 物理探査の観測対象と装置を比較し、海上音波探査の受振器を誤って説明した記述を選ぶ。

\paragraph{問われる核心} 観測量と受振器の対応

\paragraph{必要知識} ハイドロフォンの圧力応答、地上速度計、相対重力計、MT、電磁波の反射。

\paragraph{考え方（査読待ち初稿）} 測定量を対照する。重力は加速度、MTは自然電磁応答、レーダは電磁反射である。水中の流体はせん断波を伝えず、海上の受振系では音圧を検出するハイドロフォンを用いる。速度計を防水加工するという第5肢の説明が誤りである。

\paragraph{正答を選ぶ理由} 測定量を対照する。重力は加速度、MTは自然電磁応答、レーダは電磁反射である。水中の流体はせん断波を伝えず、海上の受振系では音圧を検出するハイドロフォンを用いる。速度計を防水加工するという第5肢の説明が誤りである。

\paragraph{選択肢の検討（査読待ち初稿）}\mbox{}\par

\noindent\textbf{1：適切} 磁気から鉄類の有無と概略位置を探るが、細かな形状や製品種類は判別できない。通常の磁気異常だけでは磁化・形状・深度が非一意であり、細かな製品識別までは確定できないという探査上の限界を述べる。現代の補助情報を用いる識別可能性全般を否定する命題とは解さない。\par

\noindent\textbf{2：適切} MTは自然の地磁気・地電流から深部比抵抗を探る。自然変動の水平電場・磁場の応答から電気抵抗率を推定する方法である。\par

\noindent\textbf{3：適切} 相対重力計は吊った錘のばね伸びを利用する。ばね定数と伸びの変化から重力の差を測る。絶対値を直接定める落下式とは区別する。\par

\noindent\textbf{4：適切} 地中レーダは概ね10 MHz以上の電磁波で極浅部を探る。原典が述べる通常の浅部レーダの周波数帯と原理に対応する。到達深度は周波数だけでなく導電率や減衰にも依存する。\par

\noindent\textbf{5：不適切} ストリーマケーブルには防水加工した速度計を多数入れる。水中音波を測る受振器は通常ハイドロフォンであり、圧力変動を検出する。地上用速度計の防水化と同一ではない。\par

\paragraph{再利用する考え方} 観測量と受振器の対応を、測定量・推定対象・基準面・近似条件に分けて理解する。名称だけで探査法を対応させない。

\paragraph{共通理論} \hyperref[theory:exploration:observables]{第\ref*{theory:exploration:observables}節}, \hyperref[theory:exploration:electromagnetic]{第\ref*{theory:exploration:electromagnetic}節}, \hyperref[theory:gravity:measurement]{第\ref*{theory:gravity:measurement}節}

\par 参考：\cite{usgs_marine_seismic,usgs_mt_metadata}。

\par 難易度 3（校準初稿）、優先度未評価。

\subsection{平成23年度 問32}
\label{q:h23-reg-applied-32}

\noindent\textbf{原典：} sources/H23/17\_applied\_science.pdf、PDF物理11頁。

\textbf{公式正答：} 3（single）。

\par\textit{原典本文の校正と独自解説は未作成。分類のみの登録項目。}

\subsection{平成24年度 問23}
\label{q:h24-reg-applied-23}

\noindent\textbf{原典：} sources/H24/17\_applied\_science.pdf、PDF物理7頁。

\textbf{公式正答：} 4（single）。

\par\textit{原典本文の校正と独自解説は未作成。分類のみの登録項目。}

\subsection{平成25年度 問32}
\label{q:h25-reg-applied-32}

\noindent\textbf{原典：} sources/H25/17\_applied\_science.pdf、PDF物理12頁。

\textbf{公式正答：} 3（single）。

\par\textit{原典本文の校正と独自解説は未作成。分類のみの登録項目。}

\subsection{平成26年度 問20}
\label{q:h26-reg-applied-20}

\noindent\textbf{原典：} sources/H26/17\_applied\_science.pdf、PDF物理8頁。

\textbf{公式正答：} 1（single）。

\par\textit{原典本文の校正と独自解説は未作成。分類のみの登録項目。}

\subsection{平成29年度 問34}
\label{q:h29-reg-applied-34}

\noindent\textbf{原典：} sources/H29/17\_applied\_science.pdf、PDF物理12頁。

\textbf{公式正答：} 3（single）。

\par\textit{原典本文の校正と独自解説は未作成。分類のみの登録項目。}

\subsection{平成30年度 問33}
\label{q:h30-reg-applied-33}

\noindent\textbf{原典：} sources/H30/17\_applied\_science.pdf、PDF物理11頁。

\textbf{公式正答：} 5（single）。

\par\textit{原典本文の校正と独自解説は未作成。分類のみの登録項目。}

\subsection{令和元年度 問34}
\label{q:r01-reg-applied-34}

\noindent\textbf{原典：} sources/R01/17\_applied\_science.pdf、PDF物理12頁。

\textbf{公式正答：} 4（single）。

\par\textit{原典本文の校正と独自解説は未作成。分類のみの登録項目。}

\subsection{令和元年度再試験 問32}
\label{q:r01-ret-applied-32}

\noindent\textbf{原典：} sources/R01\_retest/17\_applied\_science.pdf、PDF物理11頁。

\textbf{公式正答：} 4（single）。

\par\textit{原典本文の校正と独自解説は未作成。分類のみの登録項目。}

\subsection{令和2年度 問32}
\label{q:r02-reg-applied-32}

\noindent\textbf{原典：} sources/R02/17\_applied\_science.pdf、PDF物理11頁。

\textbf{公式正答：} 2（single）。

\par\textit{原典本文の校正と独自解説は未作成。分類のみの登録項目。}

\subsection{令和2年度 問35}
\label{q:r02-reg-applied-35}

\noindent\textbf{原典：} sources/R02/17\_applied\_science.pdf、PDF物理12頁。

\textbf{公式正答：} 2（single）。

\par\textit{原典本文の校正と独自解説は未作成。分類のみの登録項目。}

\subsection{令和3年度 問33}
\label{q:r03-reg-applied-33}

\noindent\textbf{原典：} sources/R03/17\_applied\_science.pdf、PDF物理12頁。

\textbf{公式正答：} 1（single）。

\par\textit{原典本文の校正と独自解説は未作成。分類のみの登録項目。}

\subsection{令和3年度 問34}
\label{q:r03-reg-applied-34}

\noindent\textbf{原典：} sources/R03/17\_applied\_science.pdf、PDF物理13頁。

\textbf{公式正答：} 2（single）。

\par\textit{原典本文の校正と独自解説は未作成。分類のみの登録項目。}

\subsection{令和4年度 問32}
\label{q:r04-reg-applied-32}

\noindent\textbf{原典：} sources/R04/17\_applied\_science.pdf、PDF物理11頁。

\textbf{公式正答：} 3（single）。

\par\textit{原典本文の校正と独自解説は未作成。分類のみの登録項目。}

\subsection{令和5年度 問32}
\label{q:r05-reg-applied-32}

\noindent\textbf{原典：} sources/R05/17\_applied\_science.pdf、PDF物理11頁。

\textbf{公式正答：} 4（single）。

\par\textit{原典本文の校正と独自解説は未作成。分類のみの登録項目。}

\subsection{令和6年度 問32}
\label{q:r06-reg-applied-32}

\noindent\textbf{原典：} sources/R06/17\_applied\_science.pdf、PDF物理11頁。

\textbf{公式正答：} 1（single）。

\par\textit{原典本文の校正と独自解説は未作成。分類のみの登録項目。}

\subsection{令和7年度 問32}
\label{q:r07-reg-applied-32}

\noindent\textbf{原典：} sources/R07/17\_applied\_science.pdf、PDF物理12頁。

\textbf{公式正答：} 3（single）。

\par\textit{原典本文の校正と独自解説は未作成。分類のみの登録項目。}

\section{重力・磁気探査}

\subsection{平成27年度 問22}
\label{q:h27-reg-applied-22}

\noindent\textbf{原典：} sources/H27/17\_applied\_science.pdf、PDF物理8頁。

\textbf{公式正答：} 2（single）。

\paragraph{問題の要旨} 地形、無限水平層、潮汐、ばねドリフト、緯度の重力効果から誤りを選ぶ。

\paragraph{問われる核心} 無限平板の重力と重力補正の基準

\paragraph{必要知識} ニュートン引力積分、無限平板と有限岩体の違い、地形補正の符号、潮汐、ドリフト、正規重力。

\paragraph{考え方（査読待ち初稿）} 無限水平層の上に観測点がある場合、各薄層の寄与を積分すると\(g_z=2\pi G\rho t\)である。層と測点の距離は消えるので、浅いほど大きいという第2肢が誤りである。第1肢は無限平板に対する地形差、第4肢は短時間の線形ドリフト近似として条件を明示して読む。

\paragraph{正答を選ぶ理由} 無限水平層の上に観測点がある場合、各薄層の寄与を積分すると\(g_z=2\pi G\rho t\)である。層と測点の距離は消えるので、浅いほど大きいという第2肢が誤りである。第1肢は無限平板に対する地形差、第4肢は短時間の線形ドリフト近似として条件を明示して読む。

\paragraph{選択肢の検討（査読待ち初稿）}\mbox{}\par

\noindent\textbf{1：適切} 山・谷の地形効果はともに重力を小さくする。測点高まで無限平板があるというブーゲー基準に対する地形効果として読む。測点より上の山の引力は上向き、谷は基準平板の下向き引力を欠くため、いずれも下向き重力を減らす。任意の山と谷の絶対重力を比較する命題ではない。\par

\noindent\textbf{2：不適切} 密度と厚さ一定の無限水平層は浅いほど重力が大きい。無限平板の外部重力は2πGρtで深度に依存しない。有限の岩体では距離依存があるため、無限という条件が重要である。\par

\noindent\textbf{3：適切} 月と太陽の引力による潮汐補正が必要である。月・太陽による時間変動を地下構造による異常と区別する必要がある。\par

\noindent\textbf{4：適切} ばねの時間変化を一定と近似し閉合誤差を各点へ配分する。同一点への戻り測定の閉合差から時間に比例するドリフトを近似して補正する操作を述べる。一定の伸びは測定区間のドリフトを一定率とする通常の近似として解釈し、あらゆる時間変化が必ず線形とはしない。\par

\noindent\textbf{5：適切} 絶対重力と緯度から求めた正規重力との差をとるのが緯度補正である。緯度に依存する理論基準を差し引き、基準地球の変化を除く操作である。\par

\paragraph{再利用する考え方} 無限平板の重力と重力補正の基準を、測定量・推定対象・基準面・近似条件に分けて理解する。名称だけで探査法を対応させない。

\paragraph{共通理論} \hyperref[theory:gravity:slab]{第\ref*{theory:gravity:slab}節}, \hyperref[theory:gravity:corrections]{第\ref*{theory:gravity:corrections}節}, \hyperref[theory:gravity:measurement]{第\ref*{theory:gravity:measurement}節}

\par 参考：\cite{gsi_gravity_definitions}。

\paragraph{関連問題} \hyperref[q:r07-reg-applied-33]{R07-REG-APPLIED-33}

\par 難易度 3（校準初稿）、優先度未評価。

\subsection{平成29年度 問31}
\label{q:h29-reg-applied-31}

\noindent\textbf{原典：} sources/H29/17\_applied\_science.pdf、PDF物理10頁。

\textbf{公式正答：} 5（single）。

\paragraph{問題の要旨} 重力探査で用いる補正名のうち、主に反射法地震探査に属するものを選ぶ。

\paragraph{問われる核心} 重力補正と地震走時補正の区別

\paragraph{必要知識} 運動・質量・潮汐・標高の重力効果、浅部層による走時差。

\paragraph{考え方（査読待ち初稿）} 補正の対象量を並べると、第1〜4肢は重力加速度の変化を扱う。静補正は主に地震探査の走時をそろえるもので、重力探査の標準的補正名としては該当しない。

\paragraph{正答を選ぶ理由} 補正の対象量を並べると、第1〜4肢は重力加速度の変化を扱う。静補正は主に地震探査の走時をそろえるもので、重力探査の標準的補正名としては該当しない。

\paragraph{選択肢の検討（査読待ち初稿）}\mbox{}\par

\noindent\textbf{1：適切} エトベス補正。船や航空機の運動に伴う回転系の見掛加速度の差を補正する。\par

\noindent\textbf{2：適切} ブーゲー補正。基準面と測点の間の岩石質量の引力を除く。\par

\noindent\textbf{3：適切} 潮汐補正。月・太陽による時間変動を除く。\par

\noindent\textbf{4：適切} フリーエア補正。基準面からの高さに伴う重力変化を補正する。\par

\noindent\textbf{5：不適切} 静補正。浅部低速度層や標高による地震波の走時差を、共通の基準面へ補正する地震探査の用語である。\par

\paragraph{再利用する考え方} 重力補正と地震走時補正の区別を、測定量・推定対象・基準面・近似条件に分けて理解する。名称だけで探査法を対応させない。

\paragraph{共通理論} \hyperref[theory:gravity:corrections]{第\ref*{theory:gravity:corrections}節}, \hyperref[theory:exploration:observables]{第\ref*{theory:exploration:observables}節}

\par 参考：\cite{gsi_gravity_definitions}。

\par 難易度 3（校準初稿）、優先度未評価。

\subsection{平成30年度 問32}
\label{q:h30-reg-applied-32}

\noindent\textbf{原典：} sources/H30/17\_applied\_science.pdf、PDF物理11頁。

\textbf{公式正答：} 3（single）。

\paragraph{問題の要旨} 重力・磁気探査の基準面、磁化特性、移動観測補正と、別の観測分野の語を区別する。

\paragraph{問われる核心} 重力・磁気探査の物理量と用語

\paragraph{必要知識} 等ポテンシャル面、残留磁化と誘導磁化、磁気秩序の温度依存、回転座標系。

\paragraph{考え方（査読待ち初稿）} ジオイドとエトベス補正は重力、ケーニヒスベルガー比とキュリー温度は磁気に直接関係する。プランク関数は熱平衡放射のスペクトルを記述する関数であり、この候補群では重力・磁気探査の基本量に当たらない。したがってこの肢を選ぶ。

\paragraph{正答を選ぶ理由} ジオイドとエトベス補正は重力、ケーニヒスベルガー比とキュリー温度は磁気に直接関係する。プランク関数は熱平衡放射のスペクトルを記述する関数であり、この候補群では重力・磁気探査の基本量に当たらない。したがってこの肢を選ぶ。

\paragraph{選択肢の検討（査読待ち初稿）}\mbox{}\par

\noindent\textbf{1：適切} ジオイド。基準となる重力等ポテンシャル面であり、重力場と標高の基準に関係する。\par

\noindent\textbf{2：適切} ケーニヒスベルガー比。自然残留磁化の大きさと現在の場による誘導磁化の大きさの比であり、磁気異常の解釈に関係する。\par

\noindent\textbf{3：不適切} プランク関数。熱平衡放射のスペクトルを記述する関数であり、この候補群では重力・磁気探査の基本量に当たらない。\par

\noindent\textbf{4：適切} キュリー温度。強磁性・フェリ磁性の自発磁化が消失する温度に対応し、磁性鉱物や磁気探査の解釈に関係する。\par

\noindent\textbf{5：適切} エトベス補正。移動観測系の運動に伴う見掛重力の差を補正する。\par

\paragraph{再利用する考え方} 重力・磁気探査の物理量と用語を、測定量・推定対象・基準面・近似条件に分けて理解する。名称だけで探査法を対応させない。

\paragraph{共通理論} \hyperref[theory:gravity:corrections]{第\ref*{theory:gravity:corrections}節}, \hyperref[theory:gravity:geoid]{第\ref*{theory:gravity:geoid}節}, \hyperref[theory:geomagnetism:magnetization]{第\ref*{theory:geomagnetism:magnetization}節}, \hyperref[theory:exploration:observables]{第\ref*{theory:exploration:observables}節}

\par 参考：\cite{gsi_gravity_definitions,kanaya_okuma2011}。

\paragraph{関連問題} \hyperref[q:r04-reg-applied-35]{R04-REG-APPLIED-35}

\par 難易度 3（校準初稿）、優先度未評価。

\subsection{令和4年度 問35}
\label{q:r04-reg-applied-35}

\noindent\textbf{原典：} sources/R04/17\_applied\_science.pdf、PDF物理12頁。

\textbf{公式正答：} 4（single）。

\paragraph{問題の要旨} 重力・磁気探査の基準面、磁化特性、移動観測補正と、別の観測分野の語を区別する。

\paragraph{問われる核心} 重力・磁気探査の物理量と用語

\paragraph{必要知識} 等ポテンシャル面、残留磁化と誘導磁化、磁気秩序の温度依存、回転座標系。

\paragraph{考え方（査読待ち初稿）} ジオイドとエトベス補正は重力、ケーニヒスベルガー比とキュリー温度は磁気に直接関係する。プランク関数は熱平衡放射のスペクトルを記述する関数であり、この候補群では重力・磁気探査の基本量に当たらない。したがってこの肢を選ぶ。

\paragraph{正答を選ぶ理由} ジオイドとエトベス補正は重力、ケーニヒスベルガー比とキュリー温度は磁気に直接関係する。プランク関数は熱平衡放射のスペクトルを記述する関数であり、この候補群では重力・磁気探査の基本量に当たらない。したがってこの肢を選ぶ。

\paragraph{選択肢の検討（査読待ち初稿）}\mbox{}\par

\noindent\textbf{1：適切} エトベス補正。移動観測系の運動に伴う見掛重力の差を補正する。\par

\noindent\textbf{2：適切} ジオイド。基準となる重力等ポテンシャル面であり、重力場と標高の基準に関係する。\par

\noindent\textbf{3：適切} キュリー温度。強磁性・フェリ磁性の自発磁化が消失する温度に対応し、磁性鉱物や磁気探査の解釈に関係する。\par

\noindent\textbf{4：不適切} プランク関数。熱平衡放射のスペクトルを記述する関数であり、この候補群では重力・磁気探査の基本量に当たらない。\par

\noindent\textbf{5：適切} ケーニヒスベルガー比。自然残留磁化の大きさと現在の場による誘導磁化の大きさの比であり、磁気異常の解釈に関係する。\par

\paragraph{再利用する考え方} 重力・磁気探査の物理量と用語を、測定量・推定対象・基準面・近似条件に分けて理解する。名称だけで探査法を対応させない。

\paragraph{共通理論} \hyperref[theory:gravity:corrections]{第\ref*{theory:gravity:corrections}節}, \hyperref[theory:gravity:geoid]{第\ref*{theory:gravity:geoid}節}, \hyperref[theory:geomagnetism:magnetization]{第\ref*{theory:geomagnetism:magnetization}節}, \hyperref[theory:exploration:observables]{第\ref*{theory:exploration:observables}節}

\par 参考：\cite{gsi_gravity_definitions,kanaya_okuma2011}。

\paragraph{関連問題} \hyperref[q:h30-reg-applied-32]{H30-REG-APPLIED-32}, \hyperref[q:r06-reg-applied-34]{R06-REG-APPLIED-34}

\par 難易度 3（校準初稿）、優先度未評価。

\subsection{令和6年度 問34}
\label{q:r06-reg-applied-34}

\noindent\textbf{原典：} sources/R06/17\_applied\_science.pdf、PDF物理12頁。

\textbf{公式正答：} 3（single）。

\paragraph{問題の要旨} 重力・磁気探査の基準面、磁化特性、移動観測補正と、別の観測分野の語を区別する。

\paragraph{問われる核心} 重力・磁気探査の物理量と用語

\paragraph{必要知識} 等ポテンシャル面、残留磁化と誘導磁化、磁気秩序の温度依存、回転座標系。

\paragraph{考え方（査読待ち初稿）} ジオイドとエトベス補正は重力、ケーニヒスベルガー比とキュリー温度は磁気に直接関係する。AE（アコースティック・エミッション）は微小破壊等に伴う過渡的な弾性波の放出・計測である。岩石研究にも用いるが、重力や磁化の測定とは異なる物理量を扱う。したがってこの肢を選ぶ。

\paragraph{正答を選ぶ理由} ジオイドとエトベス補正は重力、ケーニヒスベルガー比とキュリー温度は磁気に直接関係する。AE（アコースティック・エミッション）は微小破壊等に伴う過渡的な弾性波の放出・計測である。岩石研究にも用いるが、重力や磁化の測定とは異なる物理量を扱う。したがってこの肢を選ぶ。

\paragraph{選択肢の検討（査読待ち初稿）}\mbox{}\par

\noindent\textbf{1：適切} ケーニヒスベルガー比。自然残留磁化の大きさと現在の場による誘導磁化の大きさの比であり、磁気異常の解釈に関係する。\par

\noindent\textbf{2：適切} ジオイド。基準となる重力等ポテンシャル面であり、重力場と標高の基準に関係する。\par

\noindent\textbf{3：不適切} AE（アコースティック・エミッション）。微小破壊等に伴う過渡的な弾性波の放出・計測である。岩石研究にも用いるが、重力や磁化の測定とは異なる物理量を扱う。\par

\noindent\textbf{4：適切} キュリー温度。強磁性・フェリ磁性の自発磁化が消失する温度に対応し、磁性鉱物や磁気探査の解釈に関係する。\par

\noindent\textbf{5：適切} エトベス補正。移動観測系の運動に伴う見掛重力の差を補正する。\par

\paragraph{再利用する考え方} 重力・磁気探査の物理量と用語を、測定量・推定対象・基準面・近似条件に分けて理解する。名称だけで探査法を対応させない。

\paragraph{共通理論} \hyperref[theory:gravity:corrections]{第\ref*{theory:gravity:corrections}節}, \hyperref[theory:gravity:geoid]{第\ref*{theory:gravity:geoid}節}, \hyperref[theory:geomagnetism:magnetization]{第\ref*{theory:geomagnetism:magnetization}節}, \hyperref[theory:exploration:observables]{第\ref*{theory:exploration:observables}節}

\par 参考：\cite{gsi_gravity_definitions,kanaya_okuma2011,sato_lei2009_ae}。

\paragraph{関連問題} \hyperref[q:r04-reg-applied-35]{R04-REG-APPLIED-35}

\par 難易度 3（校準初稿）、優先度未評価。

\subsection{令和7年度 問33}
\label{q:r07-reg-applied-33}

\noindent\textbf{原典：} sources/R07/17\_applied\_science.pdf、PDF物理12頁。

\textbf{公式正答：} 5（single）。

\paragraph{問題の要旨} 正規重力と各補正を区別し、密度差によるブーゲー異常の符号を誤る記述を選ぶ。

\paragraph{問われる核心} 基準密度と重力異常の符号

\paragraph{必要知識} 正規重力、フリーエア・ブーゲー・エトベス補正、密度差の引力積分。

\paragraph{考え方（査読待ち初稿）} 測定重力から正規場・高さ・地表質量等の効果を取り除き、地下の密度差を評価する。同じ幾何条件なら重力の差は密度差\(\Delta\rho\)に線形であり、高密度は正の、低密度は負の異常を生む。第5肢はこの符号を逆に記述する。

\paragraph{正答を選ぶ理由} 測定重力から正規場・高さ・地表質量等の効果を取り除き、地下の密度差を評価する。同じ幾何条件なら重力の差は密度差\(\Delta\rho\)に線形であり、高密度は正の、低密度は負の異常を生む。第5肢はこの符号を逆に記述する。

\paragraph{選択肢の検討（査読待ち初稿）}\mbox{}\par

\noindent\textbf{1：適切} 正規重力は地球の質量・自転・形状から導く理論値。基準回転楕円体の外部重力ポテンシャルと自転による遠心ポテンシャルから定める。\par

\noindent\textbf{2：適切} ブーゲー補正は測点と海水準面の間の地殻物質の効果を除く。標高差だけでなく、その間の岩石質量の引力を平板等で近似して除く。\par

\noindent\textbf{3：適切} 船・航空機の移動は遠心効果を変える。地球自転と観測機の相対運動により見掛加速度が変わるためエトベス補正が必要である。\par

\noindent\textbf{4：適切} フリーエア補正は高さの違いの影響を補正する。高さによる重力の鉛直勾配の影響をそろえる操作である。\par

\noindent\textbf{5：不適切} 高密度地域のブーゲー異常は小さく、低密度地域では大きくなる。形状等を同条件にした基準密度との差では、密度の過剰は正、欠損は負の重力異常を生む。第5肢は基本的対応を逆にしている。\par

\paragraph{再利用する考え方} 基準密度と重力異常の符号を、測定量・推定対象・基準面・近似条件に分けて理解する。名称だけで探査法を対応させない。

\paragraph{共通理論} \hyperref[theory:gravity:slab]{第\ref*{theory:gravity:slab}節}, \hyperref[theory:gravity:corrections]{第\ref*{theory:gravity:corrections}節}

\par 参考：\cite{gsi_gravity_definitions}。

\paragraph{関連問題} \hyperref[q:h27-reg-applied-22]{H27-REG-APPLIED-22}

\par 難易度 3（校準初稿）、優先度未評価。

\section{電気・電磁探査}

\subsection{平成23年度 問30}
\label{q:h23-reg-applied-30}

\noindent\textbf{原典：} sources/H23/17\_applied\_science.pdf、PDF物理10頁。

\textbf{公式正答：} 2（single）。

\par\textit{原典本文の校正と独自解説は未作成。分類のみの登録項目。}

\subsection{平成24年度 問20}
\label{q:h24-reg-applied-20}

\noindent\textbf{原典：} sources/H24/17\_applied\_science.pdf、PDF物理7頁。

\textbf{公式正答：} 5（single）。

\par\textit{原典本文の校正と独自解説は未作成。分類のみの登録項目。}

\subsection{平成25年度 問34}
\label{q:h25-reg-applied-34}

\noindent\textbf{原典：} sources/H25/17\_applied\_science.pdf、PDF物理12頁。

\textbf{公式正答：} 3（single）。

\par\textit{原典本文の校正と独自解説は未作成。分類のみの登録項目。}

\subsection{平成25年度 問35}
\label{q:h25-reg-applied-35}

\noindent\textbf{原典：} sources/H25/17\_applied\_science.pdf、PDF物理13頁。

\textbf{公式正答：} 2（single）。

\par\textit{原典本文の校正と独自解説は未作成。分類のみの登録項目。}

\subsection{平成26年度 問22}
\label{q:h26-reg-applied-22}

\noindent\textbf{原典：} sources/H26/17\_applied\_science.pdf、PDF物理9頁。

\textbf{公式正答：} 3（single）。

\par\textit{原典本文の校正と独自解説は未作成。分類のみの登録項目。}

\subsection{平成27年度 問19}
\label{q:h27-reg-applied-19}

\noindent\textbf{原典：} sources/H27/17\_applied\_science.pdf、PDF物理7頁。

\textbf{公式正答：} 3（single）。

\paragraph{問題の要旨} 放射能・地温・電磁・レーダ・重力探査の推定対象を比べ、CSAMTの物性を誤る記述を選ぶ。

\paragraph{問われる核心} 電磁応答と比抵抗の対応

\paragraph{必要知識} 比抵抗と導電率、平面波インピーダンス、CSAMT遠方場、ガンマ線・温度・重力の観測。

\paragraph{考え方（査読待ち初稿）} CSAMTは人工源の遠方場を用いる電磁探査であり、電磁応答から比抵抗構造を推定する。均質・平面波近似なら見掛比抵抗は\(\rho_a=|E/H|^2/(\omega\mu)\)である。磁化率を主対象とする第3肢は物性の対応を誤っている。

\paragraph{正答を選ぶ理由} CSAMTは人工源の遠方場を用いる電磁探査であり、電磁応答から比抵抗構造を推定する。均質・平面波近似なら見掛比抵抗は\(\rho_a=|E/H|^2/(\omega\mu)\)である。磁化率を主対象とする第3肢は物性の対応を誤っている。

\paragraph{選択肢の検討（査読待ち初稿）}\mbox{}\par

\noindent\textbf{1：適切} 放射能探査は主にガンマ線強度から断層等の地質情報を得る。天然放射性核種の分布や移動の差を地質情報として利用する。ガンマ線の減衰により観測は主に表層からの寄与であり、深部核種をそのまま透視する意味ではない。\par

\noindent\textbf{2：適切} 地温探査は温泉・地熱・地下水等の調査に用いる。温度分布と熱の移流・伝導を手掛かりに熱水・地下水等を探る。\par

\noindent\textbf{3：不適切} CSAMTは遠方人工源の直交水平電場・磁場から磁化率を得る。主要な推定対象は比抵抗である。水平電場と磁場の比を磁化率そのものと取り違えている。\par

\noindent\textbf{4：適切} レーダは誘電率や導電率の違う境界の電磁反射を捉える。電気特性の差が波の伝播と反射を変えることを利用する。\par

\noindent\textbf{5：適切} 重力は密度構造を求め、標高精度が重要である。標高1 mの誤差は標準鉛直勾配で約0.3086 mGalに相当し、微小重力異常の解釈を妨げる。\par

\paragraph{再利用する考え方} 電磁応答と比抵抗の対応を、測定量・推定対象・基準面・近似条件に分けて理解する。名称だけで探査法を対応させない。

\paragraph{共通理論} \hyperref[theory:exploration:electromagnetic]{第\ref*{theory:exploration:electromagnetic}節}, \hyperref[theory:exploration:observables]{第\ref*{theory:exploration:observables}節}

\par 参考：\cite{ochi_csamt_saijo,gsi_gravity_definitions}。

\par 難易度 3（校準初稿）、優先度未評価。

\subsection{平成28年度 問19}
\label{q:h28-reg-applied-19}

\noindent\textbf{原典：} sources/H28/17\_applied\_science.pdf、PDF物理7頁。

\textbf{公式正答：} 1（single）。

\par\textit{原典本文の校正と独自解説は未作成。分類のみの登録項目。}

\subsection{平成29年度 問32}
\label{q:h29-reg-applied-32}

\noindent\textbf{原典：} sources/H29/17\_applied\_science.pdf、PDF物理11頁。

\textbf{公式正答：} 4（single）。

\par\textit{原典本文の校正と独自解説は未作成。分類のみの登録項目。}

\subsection{令和3年度 問32}
\label{q:r03-reg-applied-32}

\noindent\textbf{原典：} sources/R03/17\_applied\_science.pdf、PDF物理12頁。

\textbf{公式正答：} 4（single）。

\par\textit{原典本文の校正と独自解説は未作成。分類のみの登録項目。}

\section{ボーリング・検層}

\subsection{平成23年度 問31}
\label{q:h23-reg-applied-31}

\noindent\textbf{原典：} sources/H23/17\_applied\_science.pdf、PDF物理10頁。

\textbf{公式正答：} 2（single）。

\par\textit{原典本文の校正と独自解説は未作成。分類のみの登録項目。}

\subsection{平成24年度 問22}
\label{q:h24-reg-applied-22}

\noindent\textbf{原典：} sources/H24/17\_applied\_science.pdf、PDF物理7頁。

\textbf{公式正答：} 2（single）。

\par\textit{原典本文の校正と独自解説は未作成。分類のみの登録項目。}

\subsection{平成27年度 問18}
\label{q:h27-reg-applied-18}

\noindent\textbf{原典：} sources/H27/17\_applied\_science.pdf、PDF物理6頁。

\textbf{公式正答：} 4（single）。

\par\textit{原典本文の校正と独自解説は未作成。分類のみの登録項目。}

\subsection{令和元年度 問33}
\label{q:r01-reg-applied-33}

\noindent\textbf{原典：} sources/R01/17\_applied\_science.pdf、PDF物理12頁。

\textbf{公式正答：} 3（single）。

\par\textit{原典本文の校正と独自解説は未作成。分類のみの登録項目。}

\subsection{令和元年度再試験 問27}
\label{q:r01-ret-applied-27}

\noindent\textbf{原典：} sources/R01\_retest/17\_applied\_science.pdf、PDF物理9頁。

\textbf{公式正答：} 3（single）。

\par\textit{原典本文の校正と独自解説は未作成。分類のみの登録項目。}

\subsection{令和5年度 問34}
\label{q:r05-reg-applied-34}

\noindent\textbf{原典：} sources/R05/17\_applied\_science.pdf、PDF物理12頁。

\textbf{公式正答：} 4（single）。

\par\textit{原典本文の校正と独自解説は未作成。分類のみの登録項目。}

\subsection{令和7年度 問34}
\label{q:r07-reg-applied-34}

\noindent\textbf{原典：} sources/R07/17\_applied\_science.pdf、PDF物理13頁。

\textbf{公式正答：} 4（single）。

\par\textit{原典本文の校正と独自解説は未作成。分類のみの登録項目。}

\section{表面波・微動探査}

\subsection{令和4年度 問34}
\label{q:r04-reg-applied-34}

\noindent\textbf{原典：} sources/R04/17\_applied\_science.pdf、PDF物理12頁。

\textbf{公式正答：} 1（single）。

\par\textit{原典本文の校正と独自解説は未作成。分類のみの登録項目。}

\subsection{令和6年度 問33}
\label{q:r06-reg-applied-33}

\noindent\textbf{原典：} sources/R06/17\_applied\_science.pdf、PDF物理11頁。

\textbf{公式正答：} 2（single）。

\par\textit{原典本文の校正と独自解説は未作成。分類のみの登録項目。}

\subsection{令和7年度 問35}
\label{q:r07-reg-applied-35}

\noindent\textbf{原典：} sources/R07/17\_applied\_science.pdf、PDF物理13頁。

\textbf{公式正答：} 2（single）。

\par\textit{原典本文の校正と独自解説は未作成。分類のみの登録項目。}

\section{探査データ処理・逆解析}

\subsection{平成23年度 問28}
\label{q:h23-reg-applied-28}

\noindent\textbf{原典：} sources/H23/17\_applied\_science.pdf、PDF物理9頁。

\textbf{公式正答：} 4（single）。

\par\textit{原典本文の校正と独自解説は未作成。分類のみの登録項目。}

\subsection{平成24年度 問19}
\label{q:h24-reg-applied-19}

\noindent\textbf{原典：} sources/H24/17\_applied\_science.pdf、PDF物理7頁。

\textbf{公式正答：} 5（single）。

\par\textit{原典本文の校正と独自解説は未作成。分類のみの登録項目。}

\subsection{平成26年度 問21}
\label{q:h26-reg-applied-21}

\noindent\textbf{原典：} sources/H26/17\_applied\_science.pdf、PDF物理9頁。

\textbf{公式正答：} 3（single）。

\par\textit{原典本文の校正と独自解説は未作成。分類のみの登録項目。}

\subsection{平成27年度 問20}
\label{q:h27-reg-applied-20}

\noindent\textbf{原典：} sources/H27/17\_applied\_science.pdf、PDF物理7頁。

\textbf{公式正答：} 1（single）。

\par\textit{原典本文の校正と独自解説は未作成。分類のみの登録項目。}

\subsection{平成28年度 問20}
\label{q:h28-reg-applied-20}

\noindent\textbf{原典：} sources/H28/17\_applied\_science.pdf、PDF物理8頁。

\textbf{公式正答：} 2（single）。

\par\textit{原典本文の校正と独自解説は未作成。分類のみの登録項目。}

\subsection{平成29年度 問35}
\label{q:h29-reg-applied-35}

\noindent\textbf{原典：} sources/H29/17\_applied\_science.pdf、PDF物理12頁。

\textbf{公式正答：} 5（single）。

\par\textit{原典本文の校正と独自解説は未作成。分類のみの登録項目。}

\subsection{平成30年度 問34}
\label{q:h30-reg-applied-34}

\noindent\textbf{原典：} sources/H30/17\_applied\_science.pdf、PDF物理11頁。

\textbf{公式正答：} 1（single）。

\par\textit{原典本文の校正と独自解説は未作成。分類のみの登録項目。}

\subsection{平成30年度 問35}
\label{q:h30-reg-applied-35}

\noindent\textbf{原典：} sources/H30/17\_applied\_science.pdf、PDF物理12頁。

\textbf{公式正答：} 5（single）。

\par\textit{原典本文の校正と独自解説は未作成。分類のみの登録項目。}

\subsection{令和元年度 問35}
\label{q:r01-reg-applied-35}

\noindent\textbf{原典：} sources/R01/17\_applied\_science.pdf、PDF物理12頁。

\textbf{公式正答：} 2（single）。

\par\textit{原典本文の校正と独自解説は未作成。分類のみの登録項目。}

\subsection{令和4年度 問33}
\label{q:r04-reg-applied-33}

\noindent\textbf{原典：} sources/R04/17\_applied\_science.pdf、PDF物理11頁。

\textbf{公式正答：} 5（single）。

\par\textit{原典本文の校正と独自解説は未作成。分類のみの登録項目。}

\subsection{令和7年度 問31}
\label{q:r07-reg-applied-31}

\noindent\textbf{原典：} sources/R07/17\_applied\_science.pdf、PDF物理11頁。

\textbf{公式正答：} 1（single）。

\par\textit{原典本文の校正と独自解説は未作成。分類のみの登録項目。}

\section{地中レーダ}

\subsection{平成24年度 問21}
\label{q:h24-reg-applied-21}

\noindent\textbf{原典：} sources/H24/17\_applied\_science.pdf、PDF物理7頁。

\textbf{公式正答：} 3（single）。

\par\textit{原典本文の校正と独自解説は未作成。分類のみの登録項目。}

\subsection{令和2年度 問33}
\label{q:r02-reg-applied-33}

\noindent\textbf{原典：} sources/R02/17\_applied\_science.pdf、PDF物理12頁。

\textbf{公式正答：} 3（single）。

\par\textit{原典本文の校正と独自解説は未作成。分類のみの登録項目。}

\subsection{令和3年度 問35}
\label{q:r03-reg-applied-35}

\noindent\textbf{原典：} sources/R03/17\_applied\_science.pdf、PDF物理13頁。

\textbf{公式正答：} 4（single）。

\par\textit{原典本文の校正と独自解説は未作成。分類のみの登録項目。}

\subsection{令和5年度 問33}
\label{q:r05-reg-applied-33}

\noindent\textbf{原典：} sources/R05/17\_applied\_science.pdf、PDF物理12頁。

\textbf{公式正答：} 3（single）。

\par\textit{原典本文の校正と独自解説は未作成。分類のみの登録項目。}

\subsection{令和6年度 問35}
\label{q:r06-reg-applied-35}

\noindent\textbf{原典：} sources/R06/17\_applied\_science.pdf、PDF物理12頁。

\textbf{公式正答：} 2（single）。

\par\textit{原典本文の校正と独自解説は未作成。分類のみの登録項目。}

\section{未確定}

\subsection{平成25年度 問33}
\label{q:h25-reg-applied-33}

\noindent\textbf{原典：} sources/H25/17\_applied\_science.pdf、PDF物理12頁。

\textbf{公式正答：} 5（single）。

\par\textit{原典本文の校正と独自解説は未作成。分類のみの登録項目。}

\subsection{平成26年度 問19}
\label{q:h26-reg-applied-19}

\noindent\textbf{原典：} sources/H26/17\_applied\_science.pdf、PDF物理8頁。

\textbf{公式正答：} 5（single）。

\par\textit{原典本文の校正と独自解説は未作成。分類のみの登録項目。}

\subsection{平成26年度 問23}
\label{q:h26-reg-applied-23}

\noindent\textbf{原典：} sources/H26/17\_applied\_science.pdf、PDF物理10頁。

\textbf{公式正答：} 4（single）。

\par\textit{原典本文の校正と独自解説は未作成。分類のみの登録項目。}

\subsection{平成27年度 問21}
\label{q:h27-reg-applied-21}

\noindent\textbf{原典：} sources/H27/17\_applied\_science.pdf、PDF物理8頁。

\textbf{公式正答：} 1（single）。

\par\textit{原典本文の校正と独自解説は未作成。分類のみの登録項目。}

\subsection{平成28年度 問21}
\label{q:h28-reg-applied-21}

\noindent\textbf{原典：} sources/H28/17\_applied\_science.pdf、PDF物理8頁。

\textbf{公式正答：} 1（single）。

\par\textit{原典本文の校正と独自解説は未作成。分類のみの登録項目。}

\subsection{平成28年度 問22}
\label{q:h28-reg-applied-22}

\noindent\textbf{原典：} sources/H28/17\_applied\_science.pdf、PDF物理9頁。

\textbf{公式正答：} 1（single）。

\par\textit{原典本文の校正と独自解説は未作成。分類のみの登録項目。}

\subsection{平成28年度 問23}
\label{q:h28-reg-applied-23}

\noindent\textbf{原典：} sources/H28/17\_applied\_science.pdf、PDF物理9頁。

\textbf{公式正答：} 3（single）。

\par\textit{原典本文の校正と独自解説は未作成。分類のみの登録項目。}

\subsection{平成29年度 問33}
\label{q:h29-reg-applied-33}

\noindent\textbf{原典：} sources/H29/17\_applied\_science.pdf、PDF物理11頁。

\textbf{公式正答：} 3（single）。

\par\textit{原典本文の校正と独自解説は未作成。分類のみの登録項目。}

\subsection{令和元年度 問32}
\label{q:r01-reg-applied-32}

\noindent\textbf{原典：} sources/R01/17\_applied\_science.pdf、PDF物理11頁。

\textbf{公式正答：} 5（single）。

\par\textit{原典本文の校正と独自解説は未作成。分類のみの登録項目。}

\subsection{令和2年度 問34}
\label{q:r02-reg-applied-34}

\noindent\textbf{原典：} sources/R02/17\_applied\_science.pdf、PDF物理12頁。

\textbf{公式正答：} 4（single）。

\par\textit{原典本文の校正と独自解説は未作成。分類のみの登録項目。}

\subsection{令和5年度 問35}
\label{q:r05-reg-applied-35}

\noindent\textbf{原典：} sources/R05/17\_applied\_science.pdf、PDF物理13頁。

\textbf{公式正答：} 2（single）。

\par\textit{原典本文の校正と独自解説は未作成。分類のみの登録項目。}
